% GENERATED FILE. Edit canonical lessons, then run npm run generate:books.
% canonical-source-hash: fnv1a64:c92c61ab744e7d1e
% canonical-lessons: UR-C131-karva, UR-C131-khatta, UR-C131-namkin, UR-C131-tez, UR-C131-dhima

\chapter{Bitter, Sour, Salty, Fast, Slow}
\label{ch:ur-bitter-sour-salty-fast-slow}
\hlchaptermodality{131}

\begin{chapteropening}
\textbf{By the end of this chapter:} \emph{I can say bitter, sour, salty, fast and slow.}

Complete the last lesson of chapter 131: I can say bitter, sour, salty, fast and slow.
\end{chapteropening}

\section[\texorpdfstring{\ur{کڑوا} (kaṛvā)}{کڑوا (kaṛvā)}]{\ur{کڑوا} (kaṛvā) — bitter}

\label{lesson:UR-C131-karva}



\begin{quote}
Before the new one: say the Urdu for the skin, then the Urdu for the chest.
\end{quote}

\subsection*{You'll want to know: \ur{کڑوا}}
\textbf{\ur{کڑوا}} — \emph{kaṛvā} — \textquotedblleft{}bitter\textquotedblright{}.

\begin{cousinweb}
\textbf{\ur{کڑوا}} (\emph{kaṛvā}) is a tadbhava, a word inherited from Sanskrit and reshaped by sound change. It comes from \emph{kaṭuka}, \textquotedblleft{}pungent, bitter\textquotedblright{} (CDIAL 2641). CDIAL lists the Hindi form as \emph{kaṛuā}.
\end{cousinweb}

\begin{grammarlens}[title={What you've built}]
One more word for this chapter.
\end{grammarlens}

\subsection*{Your turn}
\begin{itemize}
\raggedright
  \item \emph{Say it:} \emph{kaṛvā}
  \item \emph{Say it:} \emph{kaṛvā}, once more
  \item \emph{Say it:} say \emph{kaṛvā}
  \item \emph{Recall:} say the Urdu for the skin, then the Urdu for the chest, then say \emph{kaṛvā} again
\end{itemize}

\begin{tcolorbox}[breakable,colback=teal!4,colframe=teal!35!black,title={Before you move on}]
What does \ur{کڑوا} mean? (\textquotedblleft{}Bitter\textquotedblright{}.) Say it once more.
\end{tcolorbox}
\section[\texorpdfstring{\ur{کھٹا} (khaṭṭā)}{کھٹا (khaṭṭā)}]{\ur{کھٹا} (khaṭṭā) — sour}

\label{lesson:UR-C131-khatta}



\begin{quote}
Before the new one: say the Urdu for the chest, then the Urdu for bitter.
\end{quote}

\subsection*{You'll want to know: \ur{کھٹا}}
\textbf{\ur{کھٹا}} — \emph{khaṭṭā} — \textquotedblleft{}sour\textquotedblright{}.

\begin{cousinweb}
\textbf{\ur{کھٹا}} (\emph{khaṭṭā}) goes back to Sanskrit \emph{khaṭṭa}, \textquotedblleft{}sour\textquotedblright{}, a word known only from the old word-lists (CDIAL 3777). CDIAL marks the Hindi form as taken over from Punjabi, and asks whether \emph{khaṭṭa} is connected with \emph{kaṭu}, \textquotedblleft{}bitter\textquotedblright{}.
\end{cousinweb}

\begin{grammarlens}[title={What you've built}]
One more word for this chapter.
\end{grammarlens}

\subsection*{Your turn}
\begin{itemize}
\raggedright
  \item \emph{Say it:} \emph{khaṭṭā}
  \item \emph{Say it:} \emph{khaṭṭā}, once more
  \item \emph{Say it:} say \emph{khaṭṭā}
  \item \emph{Recall:} say the Urdu for the chest, then the Urdu for bitter, then say \emph{khaṭṭā} again
\end{itemize}

\begin{tcolorbox}[breakable,colback=teal!4,colframe=teal!35!black,title={Before you move on}]
What does \ur{کھٹا} mean? (\textquotedblleft{}Sour\textquotedblright{}.) Say it once more.
\end{tcolorbox}
\section[\texorpdfstring{\ur{نمکین} (namkīn)}{نمکین (namkīn)}]{\ur{نمکین} (namkīn) — salty}

\label{lesson:UR-C131-namkin}



\begin{quote}
Before the new one: say the Urdu for bitter, then the Urdu for sour.
\end{quote}

\subsection*{You'll want to know: \ur{نمکین}}
\textbf{\ur{نمکین}} — \emph{namkīn} — \textquotedblleft{}salty\textquotedblright{}.

\begin{cousinweb}
\textbf{\ur{نمکین}} (\emph{namkīn}) is a loan, usually traced to Persian \emph{namakīn}, \textquotedblleft{}salty\textquotedblright{}. It is built on \emph{namak}, \textquotedblleft{}salt\textquotedblright{}.
\end{cousinweb}

\begin{grammarlens}[title={What you've built}]
One more word for this chapter.
\end{grammarlens}

\subsection*{Your turn}
\begin{itemize}
\raggedright
  \item \emph{Say it:} \emph{namkīn}
  \item \emph{Say it:} \emph{namkīn}, once more
  \item \emph{Say it:} say \emph{namkīn}
  \item \emph{Recall:} say the Urdu for bitter, then the Urdu for sour, then say \emph{namkīn} again
\end{itemize}

\begin{tcolorbox}[breakable,colback=teal!4,colframe=teal!35!black,title={Before you move on}]
What does \ur{نمکین} mean? (\textquotedblleft{}Salty\textquotedblright{}.) Say it once more.
\end{tcolorbox}
\section[\texorpdfstring{\ur{تیز} (tez)}{تیز (tez)}]{\ur{تیز} (tez) — fast, sharp}

\label{lesson:UR-C131-tez}



\begin{quote}
Before the new one: say the Urdu for sour, then the Urdu for salty.
\end{quote}

\subsection*{You'll want to know: \ur{تیز}}
\textbf{\ur{تیز}} — \emph{tez} — \textquotedblleft{}fast, sharp\textquotedblright{}.

\begin{cousinweb}
\textbf{\ur{تیز}} (\emph{tez}) is a loan, usually traced to Persian \emph{tēz}, \textquotedblleft{}sharp, swift\textquotedblright{}.
\end{cousinweb}

\begin{grammarlens}[title={What you've built}]
One more word for this chapter.
\end{grammarlens}

\subsection*{Your turn}
\begin{itemize}
\raggedright
  \item \emph{Say it:} \emph{tez}
  \item \emph{Say it:} \emph{tez}, once more
  \item \emph{Say it:} say \emph{tez}
  \item \emph{Recall:} say the Urdu for sour, then the Urdu for salty, then say \emph{tez} again
\end{itemize}

\begin{tcolorbox}[breakable,colback=teal!4,colframe=teal!35!black,title={Before you move on}]
What does \ur{تیز} mean? (\textquotedblleft{}Fast, sharp\textquotedblright{}.) Say it once more.
\end{tcolorbox}
\section[\texorpdfstring{\ur{دھیما} (dhīmā)}{دھیما (dhīmā)}]{\ur{دھیما} (dhīmā) — slow, quiet}

\label{lesson:UR-C131-dhima}



\begin{quote}
Before the new one: say the Urdu for salty, then the Urdu for fast.
\end{quote}

\subsection*{You'll want to know: \ur{دھیما}}
\textbf{\ur{دھیما}} — \emph{dhīmā} — \textquotedblleft{}slow, quiet\textquotedblright{}.

\begin{cousinweb}
\textbf{\ur{دھیما}} (\emph{dhīmā}) is a tadbhava, a word inherited and reshaped by sound change. It comes from \emph{*dhīmma}, \textquotedblleft{}slow\textquotedblright{}, a form CDIAL reconstructs from its descendants (CDIAL 6814). CDIAL asks whether it is related to \emph{*ḍhīmma}, \textquotedblleft{}lump\textquotedblright{}.
\end{cousinweb}

\begin{grammarlens}[title={What you've built}]
One more word for this chapter.
\end{grammarlens}

\subsection*{Your turn}
\begin{itemize}
\raggedright
  \item \emph{Say it:} \emph{dhīmā}
  \item \emph{Say it:} \emph{dhīmā}, once more
  \item \emph{Say it:} say \emph{dhīmā}
  \item \emph{Recall:} say the Urdu for salty, then the Urdu for fast, then say \emph{dhīmā} again
\end{itemize}

\begin{tcolorbox}[breakable,colback=teal!4,colframe=teal!35!black,title={Before you move on}]
What does \ur{دھیما} mean? (\textquotedblleft{}Slow, quiet\textquotedblright{}.) Say it once more.
\end{tcolorbox}
